
% ============================================================================
\clearpage
\section*{v4.9.0-r82 Source-Only: The Higgs Sector of W2, the Stability of the Matter--Tachyon Cancellation, and a Torsion Pairing Channel}
\addcontentsline{toc}{section}{v4.9.0-r82 Source-Only: The Higgs Sector of W2, the Stability of the Matter--Tachyon Cancellation, and a Torsion Pairing Channel}
\label{sec:v490-r82-higgs-stability}
% ============================================================================

\paragraph{Purpose.}
Layer v4.9.0-r81 set up the joint action W1 and found that a vacuum term needs tachyonic modes and that its size is not fixed
(\S\ref{sec:v490-r81-joint-action-w1}).  On 5 November 2026 the author asked to continue with the Higgs sector.  This layer adds it
(action W2), records the bookkeeping of the vacuum energy in the matter--tachyon pairing, two pre-registered stability tests (H74,
H75, both not blind), and an exploration of the pairing channel of the torsion interaction of W1.  The notes are
\texttt{notes/\allowbreak research/\allowbreak v4\_10\_0\_r82/}.  Earlier layers are not rewritten.

\noindent\textbf{Source-only status.}
This v4.9.0-r82 layer is source-only.  It does not build a PDF, does not assign a new DOI,
and does not authorize a public release.  The v4.9.0 Version DOI
\href{https://doi.org/10.5281/zenodo.23534954}{10.5281/zenodo.23534954} and the Concept DOI
\href{https://doi.org/10.5281/zenodo.19210728}{10.5281/zenodo.19210728} are unchanged by it.

\subsection*{v4.9.0-r82 results}
\label{subsec:v490-r82-results}

\paragraph{Theoremlet v4.9.0-r82.1 (W2 and its vacuum energy).}
Exploration, analytic.  W2 is W1 plus a Higgs doublet with \(V=-\mu^2|\Phi|^2+\lambda|\Phi|^4\) and Yukawa couplings through the
connectors of r44.1(b); the sign of \(-\mu^2\) is read as the tachyonic transition connector of r81.9(c), not derived.
\begin{enumerate}[leftmargin=2em]
\item[(a)] Alone, W2 has the vacuum energy \(V_{\min}=-\mu^4/(4\lambda)\), about \((100\,\mathrm{GeV})^4\), the known excess of about \(10^{55}\).
\item[(b)] With the tachyon image of every field and the sign F3, the whole energy--momentum tensor of the image changes sign,
including its vacuum energy.  With weights \(a\) (M) and \(b\) (T) the vacuum energy is \((a-b)V_{\min}\): exactly zero at photon
weights for every \(\mu,\lambda\); \(V_{\min}/\gamma\) at matter weights.  This is the structure of the energy parity of Kaplan and
Sundrum (2005) and the reading ``\(\Lambda=0\) in the photon frame'' of Theorem~\ref{thm:dark-energy}'s status note.
\item[(c)] A time offset between the condensates of M and T leaves the cancellation of a formed condensate exact; a residual
arises only transiently during the transitions.
\end{enumerate}

\paragraph{Theoremlet v4.9.0-r82.2 (H74: direct coupling destroys the exact cancellation).}
Pre-registered (commit \texttt{b24bb08}, not blind).  A normal oscillator coupled linearly (\(\varepsilon\phi\chi\)) to one with the sign
F3 has \(\Omega^2=\tfrac12(\omega_1^2+\omega_2^2)\pm\sqrt{(\omega_1^2-\omega_2^2)^2/4-\varepsilon^2}\).
\begin{enumerate}[leftmargin=2em]
\item[(a)] Equal spectra, as required for exact cancellation, grow for every \(\varepsilon>0\) with rate \(\varepsilon/(2\omega)\) (measured
\(5.011\times10^{-4}\), \(4.985\times10^{-3}\), \(0.0498\) against \(5\times10^{-4}\), \(5\times10^{-3}\), \(0.05\)).
\item[(b)] Detuned spectra are bounded for \(\varepsilon<|\omega_1^2-\omega_2^2|/2\) and grow above (rate \(0.01594\) against \(0.01618\)).
\item[(c)] Two normal oscillators stay bounded (control).
\end{enumerate}
Run 1 with the pre-registered time \(T=400\) left three small-coupling cases undecided (planning error: growth by \(10^3\) needs
\(T\sim14/\varepsilon\)); run 2 with adapted time confirms them.  \emph{Reading.}  Exact cancellation and stability exclude each other
under direct coupling; stability needs detuning, which leaves a residual \(\gtrsim2\varepsilon|V_{\min}|/\omega^2\).

\paragraph{Theoremlet v4.9.0-r82.3 (H75: gravity as the only coupling).}
Pre-registered (commit \texttt{6aa8477}, not blind).  Newtonian two-fluid perturbations, matter with density \(\rho_A\) and the tachyon
sector with negative inertial and gravitational mass \(-\rho_B\); both fall alike, \(\ddot x=4\pi G N x\),
\(N=\bigl(\begin{smallmatrix}\rho_A&-\rho_B\\\rho_A&-\rho_B\end{smallmatrix}\bigr)\).
\begin{enumerate}[leftmargin=2em]
\item[(a)] At photon weights \(N^2=0\): no exponential growth; the difference of the sectors stays constant and the common mode grows
as \(t^2\) (exactly \(\delta_0(1+t^2)\)).  The pre-registered bound ``at most linear'' was wrong for the common mode.
\item[(b)] At matter weights one Jeans mode with rate \(\sqrt{4\pi G\rho_A(1-\bsym^2)}\) (exact); two normal fluids give
\(\sqrt{4\pi G\rho_A(1+\bsym^2)}\).
\item[(c)] A pair \(\pm m\) accelerates jointly at constant distance and constant acceleration with conserved energy and momentum
(Bondi's runaway), confirmed exactly.
\end{enumerate}
\emph{Reading.}  Gravity as the only coupling has no resonance; the price is a joint runaway.  The quantum decay of the vacuum into
pairs of positive and negative energy is not covered.

\paragraph{Theoremlet v4.9.0-r82.4 (a pairing channel of the torsion interaction).}
Exploration.  The Fierz identity \(\bar\sigma^\mu_{ab}\bar\sigma_{\mu\,cd}=2\varepsilon_{ac}\varepsilon_{bd}\) gives
\(j_1\cdot j_2=2|\Phi|^2\) with the Lorentz-scalar pair \(\Phi=\psi_2^{\mathsf T}\varepsilon\psi_1\).  For anticommuting fields the
torsion interaction \(-\tfrac3{16}\kappa\,j\cdot j\) of r81.1(b) is attractive in this channel.  Its strength \(\tfrac34\kappa=6\pi/M_P^2\)
reaches the critical value of a condensate only for a cutoff \(\gtrsim1.5\,M_P\).  The only gauge-invariant channel is \(\nu^c\nu^c\)
(r44.1(d)); a condensate there would give a Majorana mass near \(M_P\) and, through the seesaw, \(m_\nu\sim10^{-6}\,\mathrm{eV}\).  The
electroweak Higgs is not obtained this way.

\paragraph{Proof and measurement sources.}
Scripts \texttt{k188}, \texttt{k188b}, \texttt{k189}, \texttt{k190} and the notes in \texttt{notes/\allowbreak research/\allowbreak v4\_10\_0\_r82/};
pre-registrations and results H74, H75 in \texttt{notes/\allowbreak research/\allowbreak predictions/}.  \(\square\)

\paragraph{Consequence.}
The structure supplies mechanisms (exact cancellation of the vacuum energy at photon weights, stability only for weak, gravitational
coupling, an attractive pairing channel), but no scale: the only scale W1 and W2 bring with them is the Planck scale.  OP5 is not
closed.

\subsection*{v4.9.0-r82 status in the three frames}
\label{subsec:v490-r82-frames}

\begin{itemize}[leftmargin=2em]
\item W2 is written in frame M; the tachyon sector enters through the declared continuation and the sign F3, the photon frame as the
equal-weight mixture.  Vacuum energies of M and T are compared only after this map.
\item r82.3 treats the tachyon sector as negative inertial and gravitational mass in frame M's Newtonian limit; this is the F3 reading,
not a frame-free statement.
\end{itemize}

\subsection*{v4.9.0-r82 register entries}
\label{subsec:v490-r82-register}

\begin{center}
\small
\begin{tabular}{@{}>{\raggedright\arraybackslash}p{0.34\textwidth}>{\raggedright\arraybackslash}p{0.58\textwidth}@{}}
\toprule
statement & status after v4.9.0-r82\\
\midrule
Higgs vacuum energy in the M--T pairing & cancels exactly at photon weights; \(V_{\min}/\gamma\) at matter weights\\
exact cancellation under direct coupling & unstable (H74)\\
exact cancellation, gravity only & no exponential instability, joint runaway (H75)\\
torsion pairing channel & attractive, Planck scale; \(\nu^c\) Majorana channel only gauge-invariant one\\
\bottomrule
\end{tabular}
\end{center}

\subsection*{v4.9.0-r82 ledger and guard update}
\label{subsec:v490-r82-ledger-guard}

\[
R_{\rm chk}^{v490\text{-}r82}=(1,0,0)
\]
for (stability of the cancellation decided in the classical linear and Newtonian cases; scale not fixed; OP5 not closed).

The layer spends no guard:
\[
R_{\rm glob}^{(v490\text{-}r82,+)}=(0,1,1,1,1),
\qquad
\Gamma_{\rm guard}^{v490\text{-}r82}=(1,1,1,1,1),
\qquad
\beta_4=1/256\ \text{unspent}.
\]
It does not imply exact finite \(C5\) zero closure, exact all-readout terminality,
\(C6/T2\), a global photon-clock/RH theorem, or global all-readout closure.
